\section{System Architecture}
\label{sec:architecture}

CMatrix is designed as a stateful, hierarchical multi-agent system. In this section, we describe the orchestration logic and the specialized agent hierarchy.

\subsection{Hierarchical Master-Worker Design}
The architecture of CMatrix is centered around a \textit{Master Orchestrator} that manages the global execution state and delegates domain-specific tasks to a set of specialized \textit{Worker Agents}. This decoupling allows each agent to maintain its own specialized prompt templates and toolsets while the Master focuses on high-level strategy and safety governance.

\begin{figure*}[t]
\centering
\input{assets/architecture.tex}
\caption{CMatrix System Architecture showing the orchestration loop, reasoning nodes, and the integrated HITL safety subsystem.}
\label{fig:architecture}
\end{figure*}

\subsection{Specialized Agent Layer}
CMatrix incorporates seven specialized worker agents, each mapped to a specific phase of the penetration testing lifecycle:
\begin{enumerate}
    \item \textbf{Network Agent}: Responsible for port scanning, service enumeration, and protocol analysis.
    \item \textbf{Web Agent}: Focuses on HTTP/HTTPS testing, header analysis, and SSL/TLS verification.
    \item \textbf{Auth Agent}: Specialized in credential testing, session management, and brute-force detection.
    \item \textbf{Vuln Intel Agent}: Interacts with external CVE databases and threat intelligence feeds.
    \item \textbf{Config Agent}: Audits system hardening and cloud configuration (AWS/Azure/GCP).
    \item \textbf{API Agent}: Dedicated to REST/GraphQL endpoint discovery and vulnerability testing.
    \item \textbf{Command Agent}: A low-level executor for system shell operations, subject to the highest safety constraints.
\end{enumerate}

\subsection{Stateful Orchestration via LangGraph}
Unlike traditional linear agents, CMatrix utilizes a cyclic state graph implemented in LangGraph. The global \textit{AgentState} $\mathcal{S}$ is a 17-field tuple that persists the core engagement context, including: (1) Conversation History, (2) Tool Execution Logs, (3) Discovered Findings, (4) Cached Credentials, (5) Pending Approvals, (6) Strategic Objectives, and (7) Multi-Agent Delegation Metadata. The transition function $\delta$ manages the flow between nodes while ensuring every state change is atomically persisted.

\section{Hierarchical HITL Safety Framework}
\label{sec:safety_framework}

The core novelty of CMatrix is its ability to enforce mandatory human oversight without disrupting the continuity of complex reasoning chains.

\subsection{The Safety Gate Mechanism}
The safety subsystem is implemented as a dedicated node in the LangGraph, termed the \textit{Approval Gate}. Every proposed action $a_t$ by any agent is intercepted by this gate before execution. The gate evaluates the action against the risk taxonomy defined in Eq. \ref{eq:safety_gate}.

\input{assets/math_formalization.tex}

\subsection{Asynchronous Interruption and Resumption}
When an action is classified as \textit{Suspend}, the orchestrator invokes a \textit{checkpoint interrupt}. The current state $\mathcal{S}_t$ is serialized and stored in a PostgreSQL database using the \textit{PostgresSaver} backend. The execution loop is then terminated, and a notification is sent to the human operator $\mathcal{H}$.

Once $\mathcal{H}$ provides an approval or rejection via the API, CMatrix retrieves the state $\mathcal{S}_t$ from the checkpoint store and resumes the execution from the exact point of interruption. Because the full reasoning history is persisted, the LLM agent continues its task without the "context loss" or "forgetting" typical of stateless agents. This mechanism guarantees Safety Property $P_4$ (Context Preservation).

\subsection{Risk Taxonomy}
Actions are categorized based on their potential for operational impact. Low-risk actions (e.g., \texttt{ping}, \texttt{curl} to public URLs) are auto-executed. High-risk actions (e.g., \texttt{nmap -p-}, exploit execution) require HITL approval. Critical actions (e.g., \texttt{rm -rf}, \texttt{killall}) No prohibited command bypassed the auto-reject registry ($P_1$) or the HITL semantic gate ($P_2$). Crucially, the system maintained a \textbf{False Suspend Rate of $<$5\%} for benign reconnaissance tools, ensuring that operational flow is not impeded by excessive human intervention.
